% !PRINTED 389
\weberpage{389}

Ist aber $\sigma$ von 1 verschieden, so kann man ein Ziffernpaar $\alpha, \beta$, beide unter den $1,2,3,4$ oder beide unter den $5,6,7,8$, so wählen, dass die Transposition
\[
\sigma\,(\alpha,\beta)\,\sigma^{-1} = (\alpha', \beta')
\]
von $(\alpha, \beta)$ verschieden ist. Denn nimmt man zunächst für $\alpha$ eine durch $\sigma^{-1}$ veränderte Ziffer, so ist $\alpha'$ von $\alpha$ verschieden (Bd.\,I, \S.\,154, 6.), wählt man dann für $\beta$ eine Ziffer, die durch $\sigma^{-1}$ nicht in $\alpha$ übergeht und die mit $\alpha$ zugleich in der ersten oder in der zweiten Hälfte der acht Ziffern vorkommt, die gewiss immer existirt, weil nur eine Ziffer durch $\sigma^{-1}$ in $\alpha$ übergeht, so ist $(\alpha', \beta')$ von $(\alpha, \beta)$ verschieden. Dann folgt aber, dass in der Gruppe $Q$ das Element $(\alpha,\beta)\,\sigma\,\Pi_{1,2,3,4}\,(\alpha,\beta)$, und folglich auch [nach \S.\,102, (8), (9)]
\[
(\alpha,\beta)\,\sigma\,\Pi_{1,2,3,4}\,(\alpha,\beta)\,\Pi_{1,2,3,4}\,\sigma^{-1} = (\alpha,\beta)\,\sigma\,(\alpha,\beta)\,\sigma^{-1} = (\alpha,\beta)(\alpha',\beta')
\]
vorkommt, und dies ist eine von der Identität verschiedene Permutation aus $S$. Damit ist also bewiesen:

3.\ \gesperrt{Die Gruppe $P$ der Doppeltangentengleichung ist einfach.}

Als specielle Anwendung können wir hervorheben, dass die Gruppe $P$ unter den 28 Wurzeln der Doppeltangentengleichung nur Permutationen der ersten Art bewirken kann, und dass \gesperrt{folglich die Discriminante dieser Gleichung ein Quadrat} ist.

Bezeichnen wir nämlich für den Augenblick mit $G$ die Gruppe aller Permutationen der 28 Wurzeln und mit $A$ die darin als Normaltheiler enthaltene alternirende Gruppe, so ist der grösste gemeinschaftliche Theiler von $A$ und $P$ ein Normaltheiler von $P$ und muss also mit $P$ identisch sein; d.\,h. $P$ ist in $A$ enthalten.

Wenn man die Gruppe der Permutationen aufsucht, die eine der Wurzeln der Doppeltangentengleichung, etwa $[1\,2]$, ungeändert lassen, so findet man eine Gruppe, die für die übrigen 27 Wurzeln noch transitiv ist\footnote{Diese Gruppe ist nach einem geometrischen Satze von \gesperrt{Geiser} isomorph mit der Gruppe der Gleichung 27\textsuperscript{sten} Grades, von der die Lösung des Problems der 27 Geraden auf einer Fläche dritter Ordnung abhängt. (Mathem. Annalen, Bd.\,I.).}.
